\chapter{Keys, Signatures, and Authorization}

\section{Purpose}

The previous chapters built the state-machine frame, identified assets
and authority, and explained how transactions carry inputs into
execution. This chapter studies the cryptographic mechanism most often
used to prove authority:

\begin{quote}
A digital signature proves that a signing authority authorized a
specific message under a specific verification rule.
\end{quote}

That is all it proves.

It does not prove the signer understood the message. It does not prove
the transaction was safe. It does not prove the UI was honest. It does
not prove the key was not stolen. It does not prove the signature is
valid in only one domain unless the signed message binds that domain. It
does not prove the same authority should be allowed to perform every
action associated with an address.

Digital signatures are necessary for blockchain authorization, but they
are not a complete authorization model.

This chapter's central discipline is:

\begin{lstlisting}
Never ask only "is the signature valid?"
Ask:
  valid under which public key?
  over which exact bytes?
  for which domain?
  for which operation?
  at which time?
  with which replay protection?
  causing which state transition?
\end{lstlisting}

In blockchain systems, private keys are not merely credentials. They are
portable authority. A compromised key can move assets, approve spenders,
vote, upgrade contracts, operate validators, sign oracle updates, relay
bridge messages, and authorize withdrawals. A misunderstood signature
can create the same practical result as a stolen key.

\section{Learning Objectives}

By the end of this chapter, you should be able to:

\begin{itemize}
\tightlist
\item
  Explain private keys, public keys, addresses, seed phrases, derivation
  paths, and wallet accounts without conflating them.
\item
  Compare ECDSA, EdDSA, BLS, and Schnorr at the level needed for
  blockchain security review.
\item
  Explain why nonce generation, malleability rules, subgroup checks,
  rogue-key resistance, and domain separation matter.
\item
  Distinguish transaction signing from message signing.
\item
  Explain why EIP-712-style typed data improves but does not solve
  signing safety.
\item
  Compare on-chain multisig, threshold signatures, MPC signing, and
  hardware signing.
\item
  Identify authorization failures involving replay, malleability,
  confused signing, stale approvals, key compromise, and wrong signer
  assumptions.
\item
  Write an authorization model for a protocol action.
\end{itemize}

\section{The Core Model}

A signature system has three basic operations:

\begin{lstlisting}
keygen:
  private_key, public_key

sign:
  signature = Sign(private_key, message)

verify:
  Verify(public_key, message, signature) -> valid or invalid
\end{lstlisting}

NIST FIPS 186-5 describes digital signatures as mechanisms that verify
the identity of a signatory and the integrity of signed data; signature
generation uses a private key, while verification uses the corresponding
public key.\footnote{NIST, ``FIPS 186-5: Digital Signature Standard'',
  https://nvlpubs.nist.gov/nistpubs/FIPS/NIST.FIPS.186-5.pdf.}

In blockchain systems, this becomes:

\begin{lstlisting}
private key:
  authority material

public key:
  verification material

address:
  usually a shorter identifier derived from, associated with, or assigned to a public key,
  account, script, object, or contract

signature:
  evidence that the signing authority approved specific bytes

authorization:
  the protocol decision to allow a state transition because verification and policy checks pass
\end{lstlisting}

The last line is the important one.


A minimal authorization-flow map should make the signed payload explicit:

\begin{lstlisting}
message:
  action
  parameters
  signer
  verifying contract or program
  chain ID or domain
  nonce or sequence number
  deadline or expiry
  intended recipient or beneficiary

verification:
  recover or verify signer
  check domain
  check nonce unused
  check deadline
  check signer has authority for this action in current state
  apply transition only if all checks pass
\end{lstlisting}

For review work, draw the flow before trusting the signature:

\begin{lstlisting}
user or signer
  -> signed payload
  -> verifier
  -> domain and replay checks
  -> authority policy
  -> state transition
\end{lstlisting}

Signature verification is usually only one part of authorization:

\begin{lstlisting}
require(signature_valid)
require(nonce_unused)
require(deadline_not_expired)
require(chain_id == expected)
require(verifying_contract == this)
require(action_allowed_for_signer)
require(current_state_still_matches_signed_assumptions)
\end{lstlisting}

If any of these checks is missing, a valid signature can authorize the
wrong transition.

\section{Private Keys, Public Keys, and Address Derivation}

\subsection{Private Keys Are Authority Material}

A private key is a secret value used to generate signatures. Whoever can
use the private key can usually exercise the authority attached to the
corresponding account, address, validator, role, or signer policy.

For an externally owned Ethereum account, ethereum.org states that the
private key is used to sign transactions and grants custody over funds
associated with the account.\footnote{ethereum.org, ``Ethereum
  Accounts'', https://ethereum.org/developers/docs/accounts/.} More
generally:

\begin{lstlisting}
If the protocol accepts signatures from key K as authority for action A,
then control of K is control of A.
\end{lstlisting}

This statement is deliberately broader than ``control of funds.'' A key
may control:

\begin{itemize}
\tightlist
\item
  A user account.
\item
  A token allowance.
\item
  A governance vote.
\item
  A validator duty.
\item
  A bridge validator role.
\item
  An oracle publisher role.
\item
  A deployment account.
\item
  A proxy admin.
\item
  A backend withdrawal signer.
\item
  A session key.
\item
  A recovery guardian.
\item
  A threshold-signing share.
\end{itemize}

Key security must be evaluated according to what the key can authorize,
not what the UI labels it.

\subsection{Public Keys Verify; Addresses Identify}

A public key is derived from a private key and used to verify
signatures. An address is usually a compact identifier derived from a
public key, script, account, object, or contract.

These are not the same object.

Ethereum:

\begin{lstlisting}
private key -> public key -> Keccak-256(public key) -> last 20 bytes -> address
\end{lstlisting}

Ethereum.org describes public address derivation as taking the last 20
bytes of the Keccak-256 hash of the public key and prefixing it with
\passthrough{\lstinline!0x!}.\footnote{ethereum.org, ``Ethereum
  Accounts'', https://ethereum.org/developers/docs/accounts/.}

Bitcoin-style systems may use addresses derived from hashes of public
keys or scripts. Solana account addresses are 32-byte public keys for
ordinary keypair accounts, while program-derived addresses are derived
addresses without private keys and are controlled by program logic. Some
systems use account addresses that are not direct hashes of public keys
at all.

Security consequences:

\begin{itemize}
\tightlist
\item
  An address is not always a public key.
\item
  A contract address may have no private key.
\item
  A program-derived address may be controlled by program rules, not by
  an externally held private key.
\item
  A smart account may validate signatures using arbitrary logic.
\item
  The same public key may correspond to different address formats on
  different chains.
\item
  A signature over bytes does not automatically prove control of every
  address format that can be derived from the same key material.
\end{itemize}

Do not say ``this address signed.'' Addresses do not sign. Keys or
account-validation logic sign or validate.

\subsection{Key Generation Depends on Entropy}

Private keys must be generated with sufficient unpredictability.

Bad key generation breaks everything downstream:

\begin{itemize}
\tightlist
\item
  Weak random number generators.
\item
  Predictable seeds.
\item
  Browser entropy failures.
\item
  Brain wallets from human-chosen phrases.
\item
  Reused test keys.
\item
  Deterministic derivation from public data.
\item
  Leaked seed generation logs.
\item
  Insecure hardware or firmware.
\end{itemize}

The math may be sound and the protocol may be correct, but if the
private key is guessable, the account is gone.

Security questions:

\begin{lstlisting}
Who generated the key?
What entropy source was used?
Was generation offline or online?
Was the seed ever displayed, transmitted, logged, or backed up insecurely?
Can the same seed generate keys on multiple chains?
Can a compromised derivation path expose more authority than intended?
\end{lstlisting}

Key generation is part of the security model, not a wallet UX detail.

\subsection{Seed Phrases Are Root Authority}

Many wallets derive many keys from one root secret. BIP-39 defines a
mnemonic sentence as a way to encode entropy and derive a binary seed
from that mnemonic using PBKDF2 with HMAC-SHA512.\footnote{Marek
  Palatinus et al., ``BIP-39: Mnemonic code for generating deterministic
  keys'',
  https://github.com/bitcoin/bips/blob/master/bip-0039.mediawiki.}

The security implication is severe:

\begin{lstlisting}
Compromise of the seed phrase compromises every descendant key derived from it,
unless another mechanism prevents use of those keys.
\end{lstlisting}

A seed phrase can control:

\begin{itemize}
\tightlist
\item
  Multiple accounts.
\item
  Multiple chains.
\item
  Multiple address types.
\item
  Hidden accounts at different derivation paths.
\item
  Legacy and modern wallet accounts.
\item
  Testnet and mainnet accounts.
\end{itemize}

Users often think a seed phrase backs up ``this wallet app.'' In
reality, it may back up a tree of cryptographic authority.

Seed phrase review questions:

\begin{itemize}
\tightlist
\item
  Is the mnemonic generated with enough entropy?
\item
  Is it shown only on a trusted display?
\item
  Is there an optional passphrase?
\item
  Are backups resistant to theft, loss, fire, coercion, and inheritance
  failure?
\item
  Are users told which accounts and chains the seed controls?
\item
  Can imported seeds be scanned by malicious software?
\item
  Are old unused accounts still holding allowances or roles?
\end{itemize}

A seed phrase is not a password. A password can often be reset by an
authority. A seed phrase often is the authority.

\subsection{Hierarchical Deterministic Wallets Create Authority Trees}

BIP-32 defines hierarchical deterministic wallets, where a tree of
keypairs is derived from a single seed. It supports extended private
keys and extended public keys, with hardened and non-hardened child
derivation.\footnote{Pieter Wuille, ``BIP-32: Hierarchical Deterministic
  Wallets'',
  https://github.com/bitcoin/bips/blob/master/bip-0032.mediawiki.}

This is powerful because:

\begin{itemize}
\tightlist
\item
  One backup can recover many keys.
\item
  Public child keys can be derived from an extended public key for
  non-hardened paths.
\item
  Services can generate receiving addresses without holding spending
  keys.
\item
  Accounts can be separated by purpose, coin type, account index,
  change, and address index.
\end{itemize}

It is also dangerous.

An extended private key can derive descendant private keys. An extended
public key can reveal descendant public keys for non-hardened paths. In
some derivation designs, exposure of a child private key together with a
parent extended public key can compromise the parent private key; BIP-32
discusses hardened derivation partly to reduce parent-key leakage
risk.\footnote{Pieter Wuille, ``BIP-32: Hierarchical Deterministic
  Wallets'',
  https://github.com/bitcoin/bips/blob/master/bip-0032.mediawiki.}

Security questions:

\begin{lstlisting}
Which root seed controls this account?
Which derivation path produced this key?
Is the path hardened where it needs to be?
Has an xpub been exposed?
Can the xpub reveal addresses, balances, privacy, or future receiving addresses?
Can one compromised child key endanger other branches?
\end{lstlisting}

Derivation paths are not metadata. They are part of the authority map.

\subsection{Wallets Are Interfaces to Authority}

Ethereum.org makes a useful distinction: an account is not a wallet; a
wallet is an interface or application that lets a user interact with an
account.\footnote{ethereum.org, ``Ethereum Accounts'',
  https://ethereum.org/developers/docs/accounts/.}

This distinction matters because wallets mediate:

\begin{itemize}
\tightlist
\item
  Key generation.
\item
  Key storage.
\item
  Address derivation.
\item
  Transaction construction.
\item
  Message display.
\item
  Signature requests.
\item
  Chain switching.
\item
  dApp permissions.
\item
  Transaction simulation.
\item
  Hardware-signer communication.
\item
  Recovery flows.
\end{itemize}

The wallet may be honest, compromised, malicious, confused, outdated, or
connected to a malicious dApp.

Security review must separate:

\begin{lstlisting}
Account authority:
  What the key or smart account can authorize.

Wallet behavior:
  What the interface asks the user to sign and how it explains the result.
\end{lstlisting}

A mathematically valid signature can still be a user-side security
failure if the wallet presented the action incorrectly.

\subsection{Key Rotation and Revocation Are Hard}

In many web systems, credentials can be reset. In blockchain systems,
key rotation depends on the asset model.

Easy cases:

\begin{itemize}
\tightlist
\item
  A smart contract role can be transferred to a new key.
\item
  A multisig owner can be replaced.
\item
  A smart account can update validation logic.
\item
  A validator can update withdrawal credentials or rotate operational
  keys if the protocol supports it.
\end{itemize}

Hard cases:

\begin{itemize}
\tightlist
\item
  A lost EOA key controlling funds cannot be rotated by itself.
\item
  A compromised seed with old token approvals may remain dangerous.
\item
  An immutable contract owner set to a lost key may be unrecoverable.
\item
  A bridge validator set may require governance to rotate.
\item
  A hardware wallet compromise may require sweeping all assets and
  revoking approvals.
\end{itemize}

Key lifecycle questions:

\begin{lstlisting}
How is the key generated?
How is it stored?
How is it backed up?
How is it used?
How is use monitored?
How is it rotated?
How is authority revoked?
What happens if it is lost?
What happens if it is compromised?
\end{lstlisting}

If there is no rotation path, then key compromise is not an incident. It
is final authority loss.

\section{ECDSA, EdDSA, BLS, Schnorr, and Signature Assumptions}

\subsection{Digital Signatures Provide Authenticity and Integrity}

A digital signature is meant to provide:

\begin{itemize}
\tightlist
\item
  Data integrity: the signed message was not modified.
\item
  Origin authentication: the signer possessed the private key
  corresponding to the public key.
\item
  Non-repudiation in some legal or operational contexts.
\end{itemize}

NIST FIPS 186-5 describes digital signatures as detecting unauthorized
modification and authenticating the signatory, while also warning that
security depends on protecting private keys and that conformance to a
standard does not guarantee system security.\footnote{NIST, ``FIPS
  186-5: Digital Signature Standard'',
  https://nvlpubs.nist.gov/nistpubs/FIPS/NIST.FIPS.186-5.pdf.}

For blockchain review, translate this into:

\begin{lstlisting}
The signature proves key possession for a message.
The protocol must decide what that key is allowed to do with that message.
\end{lstlisting}

\subsection{ECDSA}

ECDSA is widely used in Bitcoin and Ethereum account transactions. FIPS
186-5 specifies ECDSA and approves deterministic ECDSA as specified in
RFC 6979.\footnote{NIST, ``FIPS 186-5: Digital Signature Standard'',
  https://nvlpubs.nist.gov/nistpubs/FIPS/NIST.FIPS.186-5.pdf.}

ECDSA security depends on:

\begin{itemize}
\tightlist
\item
  The elliptic curve group, such as secp256k1 in Bitcoin and Ethereum.
\item
  Private key secrecy.
\item
  Correct public key validation.
\item
  Correct message hashing.
\item
  Correct per-signature nonce generation.
\item
  Correct signature verification.
\item
  Correct malleability handling.
\item
  Side-channel resistance.
\end{itemize}

The per-signature nonce is especially critical. If the same ECDSA nonce
is reused across two different messages, or if nonces are biased enough,
the private key can be recovered.

Security review questions:

\begin{lstlisting}
Who signs?
Which curve?
Which hash?
How is the per-message nonce generated?
Is deterministic ECDSA used correctly?
Are signatures canonicalized?
Can signatures be malleated?
Is the recovered signer checked against the expected authority?
\end{lstlisting}

\subsection{ECDSA Malleability and Low-s Rules}

ECDSA signatures are pairs usually represented as
\passthrough{\lstinline!(r, s)!} plus sometimes a recovery identifier.
For many curves, \passthrough{\lstinline!(r, n - s)!} can also verify
for the same message and public key unless verification rules reject one
form.

This is signature malleability: a valid signature can be transformed
into a different valid signature for the same message.

Bitcoin and Ethereum have both had to care about this. Ethereum's EIP-2
made high-\passthrough{\lstinline!s!} transaction signatures invalid
while leaving the \passthrough{\lstinline!ecrecover!} precompile
behavior unchanged for compatibility with existing
signatures.\footnote{Ethereum Improvement Proposals, ``EIP-2: Homestead
  Hard-fork Changes'', https://eips.ethereum.org/EIPS/eip-2.}

Security consequences:

\begin{itemize}
\tightlist
\item
  Do not use the raw signature bytes as the unique identity of an
  authorization unless malleability is handled.
\item
  Prefer consumed message hashes, nonces, or structured authorization
  IDs for replay protection.
\item
  If using \passthrough{\lstinline!ecrecover!}, enforce valid
  \passthrough{\lstinline!v!}, \passthrough{\lstinline!r!},
  \passthrough{\lstinline!s!} ranges and low-\passthrough{\lstinline!s!}
  where appropriate.
\item
  Do not assume two byte-distinct signatures imply two different
  authorizations.
\end{itemize}

The safer question is:

\begin{lstlisting}
What message or authorization is consumed?
\end{lstlisting}

Not:

\begin{lstlisting}
Have these exact signature bytes been seen before?
\end{lstlisting}

\subsection{EdDSA}

EdDSA is a digital signature scheme over Edwards curves. RFC 8032
specifies Ed25519 and Ed448.\footnote{Josefsson and Liusvaara, ``RFC
  8032: Edwards-Curve Digital Signature Algorithm'',
  https://datatracker.ietf.org/doc/html/rfc8032.} FIPS 186-5 also
specifies EdDSA.\footnote{NIST, ``FIPS 186-5: Digital Signature
  Standard'',
  https://nvlpubs.nist.gov/nistpubs/FIPS/NIST.FIPS.186-5.pdf.}

EdDSA is commonly associated with:

\begin{itemize}
\tightlist
\item
  Deterministic signing.
\item
  Compact keys and signatures.
\item
  Fast verification.
\item
  Ed25519 usage in ecosystems such as Solana.
\end{itemize}

But EdDSA is not ``safe by default'' in every integration.

Security concerns include:

\begin{itemize}
\tightlist
\item
  Public key validation.
\item
  Canonical encoding.
\item
  Cofactor and subgroup handling.
\item
  Domain separation variants.
\item
  Prehash vs pure modes.
\item
  Incorrect library APIs.
\item
  Signing-oracle misuse.
\item
  Cross-protocol key reuse.
\end{itemize}

Security review questions:

\begin{lstlisting}
Which EdDSA variant is used?
Is the message prehashed?
Is domain separation included?
Are public keys and signatures canonical?
Does the library verify according to the ecosystem's expected rules?
Can the same key sign messages in another protocol with different assumptions?
\end{lstlisting}

Do not assume ``Ed25519'' names one complete security model. The
verification rules and message domain matter.

\subsection{BLS}

BLS signatures are pairing-based signatures with aggregation properties.
The CFRG BLS draft describes BLS as deterministic, non-malleable,
aggregation-friendly, and useful where storage or bandwidth are scarce;
it also explains that aggregate signatures can authenticate multiple
messages and public keys.\footnote{Boneh et al., ``BLS Signatures''
  Internet-Draft,
  https://datatracker.ietf.org/doc/html/draft-irtf-cfrg-bls-signature.}

Ethereum proof-of-stake uses BLS keys for validators; the Ethereum
consensus specs define BLS signature helpers such as
\passthrough{\lstinline!Sign!}, \passthrough{\lstinline!Verify!},
\passthrough{\lstinline!Aggregate!},
\passthrough{\lstinline!FastAggregateVerify!},
\passthrough{\lstinline!AggregateVerify!}, and
\passthrough{\lstinline!KeyValidate!} for beacon-chain
operations.\footnote{Ethereum Consensus Specs, ``Beacon Chain: BLS
  Signatures'',
  https://ethereum.github.io/consensus-specs/specs/phase0/beacon-chain/\#bls-signatures.}

BLS is powerful because many signatures can be compressed into one
aggregate signature. That is valuable for consensus protocols where many
validators sign votes.

But aggregation introduces specific risks:

\begin{itemize}
\tightlist
\item
  Rogue-key attacks.
\item
  Missing proof-of-possession checks.
\item
  Public key validation failures.
\item
  Subgroup membership failures.
\item
  Wrong domain separation.
\item
  Aggregating signatures over the wrong message set.
\item
  Treating aggregate success as proof that all expected signers signed.
\end{itemize}

The BLS draft explicitly defines a rogue-key attack as one involving a
specially crafted public key used to forge an aggregate signature, and
specifies methods such as proof-of-possession schemes to secure against
it.\footnote{Boneh et al., ``BLS Signatures'' Internet-Draft,
  https://datatracker.ietf.org/doc/html/draft-irtf-cfrg-bls-signature.}

Security review questions:

\begin{lstlisting}
Which BLS scheme variant is used: basic, message augmentation, or proof-of-possession?
Are public keys validated?
Are subgroup checks performed?
Is the aggregate signer set explicit?
Are duplicate public keys handled?
Is the signed domain correct?
Can an aggregate signature be replayed across fork, domain, chain, or duty?
\end{lstlisting}

BLS aggregation reduces bandwidth. It does not reduce the need to know
exactly who signed exactly what.

\subsection{Schnorr Signatures}

Schnorr signatures are important in Bitcoin Taproot and in threshold
protocols. BIP-340 specifies Schnorr signatures for secp256k1, using
32-byte public keys and 64-byte signatures, and includes design choices
such as tagged hashing for domain separation.\footnote{Pieter Wuille et
  al., ``BIP-340: Schnorr Signatures for secp256k1'',
  https://github.com/bitcoin/bips/blob/master/bip-0340.mediawiki.}

Schnorr has properties that make it attractive:

\begin{itemize}
\tightlist
\item
  Simpler structure than ECDSA.
\item
  Natural linearity useful for multisignatures and threshold signatures.
\item
  Efficient verification.
\item
  Compatibility with constructions such as MuSig and FROST.
\end{itemize}

But linearity is a tool, not a free lunch. Threshold and multisignature
protocols must handle nonce commitments, rogue keys, participant
binding, aborts, and signer coordination.

RFC 9591 specifies FROST, a two-round Schnorr threshold protocol, and
states that threshold signatures require cooperation among a threshold
number of participants, with security assuming an adversary corrupts
fewer than the threshold.\footnote{Connolly et al., ``RFC 9591: The
  Flexible Round-Optimized Schnorr Threshold Protocol'',
  https://www.rfc-editor.org/rfc/rfc9591.html.}

Security review questions:

\begin{lstlisting}
Which Schnorr variant is used?
How are nonces generated and committed?
Are public keys bound to participants?
Is the signer set committed to the challenge?
Can a coordinator equivocate?
Can nonce reuse reveal a key or share?
Can an aborting signer grief liveness?
\end{lstlisting}

Schnorr makes powerful threshold designs possible. Those designs still
have protocol risk.

\subsection{Signature Assumptions Are Broader Than Algorithms}

When a protocol says ``secured by signatures,'' ask which assumptions
are included.

{\def\LTcaptype{none} % do not increment counter
\begin{longtable}[]{@{}
  >{\raggedright\arraybackslash}p{(\linewidth - 2\tabcolsep) * \real{0.5000}}
  >{\raggedright\arraybackslash}p{(\linewidth - 2\tabcolsep) * \real{0.5000}}@{}}
\toprule\noalign{}
\begin{minipage}[b]{\linewidth}\raggedright
Assumption
\end{minipage} & \begin{minipage}[b]{\linewidth}\raggedright
Failure
\end{minipage} \\
\midrule\noalign{}
\endhead
\bottomrule\noalign{}
\endlastfoot
Private keys remain secret & Attacker signs valid messages \\
Signing devices are honest & Device signs unintended payload \\
Randomness/nonces are safe & Private key or share recovery \\
Message encoding is canonical & Signature authorizes ambiguous
message \\
Domain separation is correct & Signature replay across contexts \\
Public keys are validated & Invalid-curve or subgroup attacks \\
Signature malleability is handled & Replay or duplicate authorization
confusion \\
Signer set is explicit & Aggregate/multisig result misinterpreted \\
Verification code is correct & Forged or wrong signatures accepted \\
Users understand the prompt & Confused signing and phishing \\
\end{longtable}
}

The signature algorithm is one row in the security model, not the whole
model.

\section{Message Signing, Transaction Signing, and Domain Separation}

\subsection{Transaction Signing Authorizes a Chain-Level Request}

A transaction signature usually authorizes a specific transaction
envelope.

For an Ethereum transaction, the signed data includes transaction fields
such as nonce, recipient, value, calldata, gas parameters, and
chain-specific replay-protection fields. EIP-155 binds transaction
signing to a chain ID to reduce cross-chain replay.\footnote{Vitalik
  Buterin, ``EIP-155: Simple Replay Attack Protection'',
  https://eips.ethereum.org/EIPS/eip-155.}

Transaction signing answers:

\begin{lstlisting}
Did this account authorize this transaction envelope on this chain?
\end{lstlisting}

It does not answer:

\begin{lstlisting}
Was the target contract honest?
Was calldata understandable?
Was the oracle safe?
Was the transaction ordered safely?
Did the user understand the effect?
\end{lstlisting}

Transaction signatures are necessary authority, not sufficient safety.

\subsection{Message Signing Authorizes Off-Chain or Contract-Interpreted
Meaning}

Message signing is broader and riskier. A user may sign bytes that are
not a chain transaction:

\begin{itemize}
\tightlist
\item
  Login challenge.
\item
  Permit.
\item
  Limit order.
\item
  Intent.
\item
  Bridge withdrawal.
\item
  Governance delegation.
\item
  Session-key grant.
\item
  Marketplace listing.
\item
  Account-abstraction operation.
\item
  Contract-wallet authorization.
\item
  Terms acceptance.
\end{itemize}

The message becomes dangerous when a contract or backend treats it as
authority.

Example:

\begin{lstlisting}
sign("withdraw 100")
\end{lstlisting}

Questions:

\begin{itemize}
\tightlist
\item
  Withdraw from which protocol?
\item
  Which chain?
\item
  Which contract?
\item
  Which token?
\item
  Which recipient?
\item
  Which nonce?
\item
  Which deadline?
\item
  Can it be replayed after upgrade?
\item
  Can it be submitted by anyone?
\item
  Is it valid if the user's state changed?
\end{itemize}

If the message does not bind these facts, the signature may authorize
more than intended.

\subsection{EIP-191 and the Problem of Raw Bytes}

EIP-191 defines a signed data standard for Ethereum, including an
initial \passthrough{\lstinline!0x19!} byte and version-specific data to
prevent signed data from being valid RLP-encoded Ethereum
transactions.\footnote{Ethereum Improvement Proposals, ``EIP-191: Signed
  Data Standard'', https://eips.ethereum.org/EIPS/eip-191.}

This is a domain-separation improvement over signing arbitrary raw
bytes.

But it does not solve human comprehension. A user may still be asked to
sign a hex blob or vague text, and a contract may still interpret the
signature in a dangerous way.

Security rule:

\begin{lstlisting}
A message-signing prefix can prevent one class of replay.
It does not prove the user understood the message or that the protocol scoped it correctly.
\end{lstlisting}

\subsection{EIP-712 Typed Data}

EIP-712 defines typed structured data hashing and signing for Ethereum.
It includes a domain separator and a structured message hash, and its
specification explicitly states that it does not include replay
protection by itself.\footnote{Ethereum Improvement Proposals,
  ``EIP-712: Typed structured data hashing and signing'',
  https://eips.ethereum.org/EIPS/eip-712.}

EIP-712 is valuable because it can bind a signature to:

\begin{itemize}
\tightlist
\item
  Domain name.
\item
  Version.
\item
  Chain ID.
\item
  Verifying contract.
\item
  Structured fields.
\item
  Human-readable typed data.
\end{itemize}

But the protocol must still include nonces, deadlines, expected state,
and correct verification logic.

Good EIP-712 signing model:

\begin{lstlisting}
Domain:
  name
  version
  chainId
  verifyingContract

Message:
  action type
  signer
  recipient
  token
  amount
  nonce
  deadline
  expected state if needed
\end{lstlisting}

Bad EIP-712 signing model:

\begin{lstlisting}
Domain looks official.
Message omits nonce or deadline.
Contract accepts it forever.
\end{lstlisting}

Typed data makes dangerous signing easier to inspect. It does not make
dangerous authorization impossible.

\subsection{Permits and Signed Approvals}

ERC-2612 adds \passthrough{\lstinline!permit!}, allowing ERC-20
approvals to be made by signature rather than by an on-chain approval
transaction.\footnote{Ethereum Improvement Proposals, ``ERC-2612: Permit
  Extension for EIP-20 Signed Approvals'',
  https://eips.ethereum.org/EIPS/eip-2612.}

This is useful:

\begin{itemize}
\tightlist
\item
  Better UX.
\item
  Gasless approvals.
\item
  Composable flows.
\item
  Signer does not need native asset for gas.
\end{itemize}

It is also dangerous:

\begin{itemize}
\tightlist
\item
  A signature can grant spending authority.
\item
  A phishing site can trick a user into signing a permit.
\item
  A permit can be submitted by someone else.
\item
  Replay protection depends on the token's nonce logic.
\item
  Deadlines must be checked.
\item
  The spender and value must be clear to the user.
\end{itemize}

Permit-style review checklist:

\begin{lstlisting}
Is owner bound?
Is spender bound?
Is token contract bound?
Is chain ID bound?
Is amount bound?
Is nonce consumed?
Is deadline checked?
Can signature malleability bypass nonce tracking?
Can the spender immediately transfer after permit?
Does the UI clearly show this is an approval?
\end{lstlisting}

A permit is not ``just a signature.'' It is delegated spending
authority.

\subsection{Contract Signatures and ERC-1271}

EOAs verify signatures through public-key cryptography. Contract
accounts can define their own signature-validation logic.

ERC-1271 standardizes a method for contracts to validate whether a
signature is valid for given data.\footnote{Ethereum Improvement
  Proposals, ``ERC-1271: Standard Signature Validation Method for
  Contracts'', https://eips.ethereum.org/EIPS/eip-1271.}

This enables:

\begin{itemize}
\tightlist
\item
  Smart contract wallets.
\item
  Multisig wallets.
\item
  Session-key schemes.
\item
  Social recovery.
\item
  Account abstraction.
\item
  Policies richer than one ECDSA key.
\end{itemize}

It also creates review questions:

\begin{itemize}
\tightlist
\item
  Is signature validity time-dependent?
\item
  Can the contract's validation logic change after signing?
\item
  Does the signature remain valid after owner rotation?
\item
  Is the validation result replayable across contracts?
\item
  Does the verifier call \passthrough{\lstinline!isValidSignature!} on
  the correct contract?
\item
  Can a malicious contract return valid for arbitrary messages?
\item
  Is a signature from a contract wallet valid only under current state,
  or historical state?
\end{itemize}

For EOAs, signature validity is mostly cryptographic. For contract
accounts, signature validity is code and state.

\subsection{Domain Separation Checklist}

Any signature that authorizes a protocol action should bind:

\begin{lstlisting}
Protocol:
  name and version

Domain:
  chain ID, network, rollup, appchain, or source/destination domain

Verifier:
  contract address, program ID, module, or backend verifier identity

Action:
  function or operation type

Actor:
  signer, owner, delegate, relayer, or account

Asset:
  token, object, vault share, bridge message, or position

Parameters:
  amount, recipient, market, order terms, or proposal data

Replay guard:
  nonce, sequence, nullifier, order ID, or consumed-message ID

Time bound:
  deadline, expiry, valid-after, valid-before, block, slot, or epoch

State expectation:
  optional but important for price, chain state, order fill, or account configuration
\end{lstlisting}

The question is not whether the signature verifies. The question is
whether it verifies only for the intended transition.

\section{Multisig, Threshold Signatures, MPC, and Hardware Signing}

\subsection{Single-Key Control Is a Concentrated Failure Mode}

A single key is simple:

\begin{lstlisting}
one private key -> one signature -> one authority
\end{lstlisting}

It is also fragile:

\begin{itemize}
\tightlist
\item
  Key theft means full compromise.
\item
  Key loss means loss of authority.
\item
  Coercion of one person may be enough.
\item
  Malware on one device may be enough.
\item
  Operational mistakes are hard to stop.
\end{itemize}

High-value systems therefore distribute signing authority.

The main approaches are:

\begin{itemize}
\tightlist
\item
  On-chain multisig.
\item
  Native script multisig.
\item
  Threshold signatures.
\item
  MPC signing.
\item
  Hardware or HSM-backed signing.
\item
  Hybrid policies combining several of these.
\end{itemize}

Each improves some risks and introduces others.

\subsection{On-Chain Multisig}

An on-chain multisig is a contract that enforces a policy such as:

\begin{lstlisting}
execute transaction if at least m of n owners approve
\end{lstlisting}

Safe describes itself as smart accounts that allow users and
organizations to manage digital assets, with multisig-style owner and
threshold policies.\footnote{Safe, ``What is Safe?'',
  https://docs.safe.global/home/what-is-safe.}

Benefits:

\begin{itemize}
\tightlist
\item
  Clear on-chain policy.
\item
  Signer set and threshold may be observable.
\item
  Can support batching, modules, guards, spending limits, and timelocks.
\item
  Can rotate owners.
\item
  Can use contract-level validation such as ERC-1271.
\end{itemize}

Risks:

\begin{itemize}
\tightlist
\item
  Contract bugs.
\item
  Malicious or overpowered modules.
\item
  Owner compromise.
\item
  Threshold too low.
\item
  Signer unavailability.
\item
  Social engineering of signers.
\item
  Blind signing of batched calls.
\item
  Upgradeable multisig infrastructure.
\item
  Chain-specific execution risk.
\end{itemize}

On-chain multisig review questions:

\begin{lstlisting}
Who are the owners?
What is the threshold?
Can owners be added or removed?
Who can change the threshold?
Are modules enabled?
Can modules bypass owner threshold?
Is there a timelock?
Can transactions be simulated and decoded?
Can signers verify exactly what they approve?
\end{lstlisting}

A multisig is not automatically safe. A 2-of-3 multisig where all keys
sit on the same laptop is a one-device wallet with extra ceremony.

\subsection{Native Multisig and Script Policies}

Some chains have native multisig or script-level policies.

Bitcoin scripts can encode spending conditions such as multisig and
timelocks. The authority is enforced by the script condition on a UTXO
rather than by a separate contract account.

Benefits:

\begin{itemize}
\tightlist
\item
  Policy can be native to the spend condition.
\item
  No general-purpose contract needed.
\item
  Can combine keys and time locks.
\end{itemize}

Risks:

\begin{itemize}
\tightlist
\item
  Complex scripts can be hard to recover from.
\item
  Fee and transaction-size constraints matter.
\item
  Key loss can make funds unspendable.
\item
  Policy privacy varies by construction.
\item
  Coordinating signatures can be operationally hard.
\end{itemize}

The review question is still:

\begin{lstlisting}
Which spending paths exist, and who can satisfy each one?
\end{lstlisting}

\subsection{Threshold Signatures}

Threshold signatures distribute one signing key among multiple
participants. A threshold number cooperate to produce one signature that
verifies under one public key.

FROST, for example, lets a threshold number of participants cooperate to
compute a Schnorr signature, improving distribution of trust and
redundancy.\footnote{Connolly et al., ``RFC 9591: The Flexible
  Round-Optimized Schnorr Threshold Protocol'',
  https://www.rfc-editor.org/rfc/rfc9591.html.}

Benefits:

\begin{itemize}
\tightlist
\item
  On-chain output may look like one normal signature.
\item
  Lower fees and better privacy than some on-chain multisig schemes.
\item
  No single participant holds the full signing key.
\item
  Can support institutional signing without exposing a single key.
\end{itemize}

Risks:

\begin{itemize}
\tightlist
\item
  Distributed key generation complexity.
\item
  Signing-round coordination.
\item
  Nonce misuse or reuse.
\item
  Aborting signers causing liveness failure.
\item
  Coordinator equivocation.
\item
  Share backup and recovery complexity.
\item
  Difficult incident response if shares are suspected compromised.
\item
  Harder on-chain transparency: observers may not know the real signing
  policy.
\end{itemize}

Threshold signature review questions:

\begin{lstlisting}
How was the key generated?
Was there a trusted dealer or distributed key generation?
How many shares exist?
What threshold is required?
Who holds shares?
How are signing nonces generated and committed?
Can one participant block signing?
Can compromised shares be rotated?
Can policy changes be audited?
Does on-chain state reveal the real trust model?
\end{lstlisting}

Threshold signatures hide complexity from the chain. Reviewers must not
let the complexity hide from the threat model.

\subsection{MPC Signing}

Multiparty computation signing is a broader operational pattern where
multiple parties jointly compute a signature without reconstructing the
private key in one place.

MPC is often used by custodians, exchanges, funds, and institutional
wallets.

Benefits:

\begin{itemize}
\tightlist
\item
  No single machine holds the full key.
\item
  Signing can enforce approval workflows.
\item
  Shares can be distributed across devices, people, cloud environments,
  or vendors.
\item
  Policies can include risk engines, limits, and human approvals.
\end{itemize}

Risks:

\begin{itemize}
\tightlist
\item
  Complex implementation.
\item
  Vendor trust.
\item
  Coordinator or policy-engine compromise.
\item
  Share refresh failures.
\item
  Backup and disaster recovery failures.
\item
  Insider collusion.
\item
  Signing the wrong transaction because policy evaluation is wrong.
\item
  Poor visibility to users and auditors.
\end{itemize}

MPC should be evaluated as a system:

\begin{lstlisting}
Cryptographic protocol
+ signing policy
+ key-share storage
+ operator workflow
+ audit logging
+ recovery process
+ software supply chain
+ incident response
\end{lstlisting}

Calling something ``MPC'' is not a security proof.

\subsection{Hardware Wallets and HSMs}

Hardware signing devices and hardware security modules try to keep
private keys inside a protected boundary.

NIST FIPS 140-3 defines security requirements for cryptographic modules
across areas such as physical security, roles, authentication,
software/firmware security, sensitive security parameter management, and
mitigation of other attacks.\footnote{NIST, ``FIPS 140-3: Security
  Requirements for Cryptographic Modules'',
  https://nvlpubs.nist.gov/nistpubs/FIPS/NIST.FIPS.140-3.pdf.}

Hardware signing can reduce:

\begin{itemize}
\tightlist
\item
  Private key exposure to a general-purpose computer.
\item
  Malware extraction of raw keys.
\item
  Some remote compromise risks.
\item
  Some accidental key leakage.
\end{itemize}

It does not automatically prevent:

\begin{itemize}
\tightlist
\item
  Blind signing.
\item
  Malicious transaction construction.
\item
  Supply-chain compromise.
\item
  Firmware compromise.
\item
  Address poisoning.
\item
  Wrong-chain signing.
\item
  User approval of a malicious payload.
\item
  Coercion.
\item
  Bad backup practices.
\end{itemize}

Hardware signing review questions:

\begin{lstlisting}
Does the device display the actual action being authorized?
Can it parse the transaction type or only show hashes?
Can users verify chain, contract, token, amount, recipient, and calldata?
How are firmware updates authenticated?
How is the seed backed up?
Can the host alter what is shown?
Can the device be used in a multisig or policy-controlled flow?
\end{lstlisting}

Hardware signing protects keys. It may not protect intent.

\subsection{Choosing the Right Signing Architecture}

{\def\LTcaptype{none} % do not increment counter
\begin{longtable}[]{@{}
  >{\raggedright\arraybackslash}p{(\linewidth - 4\tabcolsep) * \real{0.3333}}
  >{\raggedright\arraybackslash}p{(\linewidth - 4\tabcolsep) * \real{0.3333}}
  >{\raggedright\arraybackslash}p{(\linewidth - 4\tabcolsep) * \real{0.3333}}@{}}
\toprule\noalign{}
\begin{minipage}[b]{\linewidth}\raggedright
Architecture
\end{minipage} & \begin{minipage}[b]{\linewidth}\raggedright
Strong against
\end{minipage} & \begin{minipage}[b]{\linewidth}\raggedright
Weak against
\end{minipage} \\
\midrule\noalign{}
\endhead
\bottomrule\noalign{}
\endlastfoot
Single EOA & Simplicity, low overhead & Theft, loss, phishing, no
policy \\
On-chain multisig & Single-key compromise, transparent policy &
Contract/module bugs, signer coordination, blind approvals \\
Script multisig & Native policy, no general contract & Policy
inflexibility, fee/size, recovery complexity \\
Threshold signature & Single-device key exposure, on-chain privacy &
Protocol complexity, hidden policy, liveness griefing \\
MPC custody & Institutional workflow, distributed key use &
Vendor/process risk, opaque failures, policy compromise \\
Hardware wallet & Key extraction from host & Blind signing, supply
chain, backup loss \\
Smart account & Rich policies, recovery, session keys & Contract bugs,
upgrade risk, validation complexity \\
\end{longtable}
}

The right question is:

\begin{lstlisting}
Which failure modes are unacceptable for this authority,
and which signing architecture actually reduces them?
\end{lstlisting}

\section{Authorization Failures: Replay, Malleability, Confused Signing, and Key Compromise}

\subsection{Signature Validity Is Not Authorization Completeness}

Authorization is a policy decision. Signature verification is one input.

Bad authorization:

\begin{lstlisting}
require(ecrecover(hash, sig) == owner)
execute()
\end{lstlisting}

Better authorization:

\begin{lstlisting}
require(ecrecover(domain_hash(action), sig) == owner)
require(nonce_unused[owner][nonce])
require(block.timestamp <= deadline)
require(action.target == expected_target)
require(action.chain_id == block.chainid)
require(action.verifying_contract == address(this))
require(action.amount <= allowed_limit)
nonce_unused[owner][nonce] = false
execute(action)
\end{lstlisting}

A valid signature over insufficient data is insufficient authorization.

\subsection{Replay Failures}

Replay occurs when an authorization can be reused outside its intended
context.

Common replay failures:

\begin{itemize}
\tightlist
\item
  Same signature used twice because no nonce is consumed.
\item
  Signature valid on multiple chains because no chain ID is bound.
\item
  Signature valid for multiple contracts because no verifying contract
  is bound.
\item
  Signature valid for multiple actions because function/type is omitted.
\item
  Signature valid after expiry because no deadline exists.
\item
  Signature valid after upgrade because version is omitted.
\item
  Bridge message valid on multiple destinations because destination is
  omitted.
\item
  Order signature partially fills repeatedly beyond intended amount.
\end{itemize}

Replay defense requires both:

\begin{lstlisting}
scope:
  signature is valid only in the intended domain

consumption:
  accepted authorization is marked used or bounded by fill accounting
\end{lstlisting}

Nonce without domain is incomplete. Domain without nonce may still
replay.

\subsection{Malleability Failures}

Malleability matters when a protocol treats signature bytes as identity.

Unsafe:

\begin{lstlisting}
used[signature] = true
\end{lstlisting}

Safer:

\begin{lstlisting}
used[signer][nonce] = true
\end{lstlisting}

or:

\begin{lstlisting}
used[authorization_hash] = true
\end{lstlisting}

where \passthrough{\lstinline!authorization\_hash!} is a canonical
structured hash of the intended action.

Malleability review questions:

\begin{lstlisting}
Can another valid signature exist for the same message?
Does verification enforce canonical signatures?
Is replay protection based on signature bytes or message identity?
Can different encodings produce the same semantic action?
\end{lstlisting}

\subsection{Confused Signing}

Confused signing happens when the signer authorizes one thing while
believing they are authorizing another.

Examples:

\begin{itemize}
\tightlist
\item
  User signs \passthrough{\lstinline!permit!} while believing they are
  logging in.
\item
  User signs unlimited token approval while believing they are swapping.
\item
  User signs a Seaport-style order without understanding consideration
  items.
\item
  Multisig signers approve a batch without decoding every call.
\item
  Hardware wallet shows only a hash.
\item
  Frontend displays token symbol from an untrusted token contract.
\item
  Wallet simulation omits delegatecall or post-signature effects.
\item
  Governance delegates vote for a proposal without inspecting executable
  calldata.
\end{itemize}

Confused signing is not solved by stronger cryptography. It is a UI,
protocol-design, and domain-separation problem.

Defenses:

\begin{itemize}
\tightlist
\item
  Typed structured signing.
\item
  Clear domain separators.
\item
  Human-readable transaction decoding.
\item
  Simulation with explicit limitations.
\item
  Spending limits.
\item
  Session-key scopes.
\item
  Timelocks.
\item
  Multisig review processes.
\item
  Hardware display of critical fields.
\item
  Warnings for approvals, permits, delegatecalls, upgrades, and unknown
  contracts.
\end{itemize}

The review question:

\begin{lstlisting}
Could a reasonable signer believe this signature is harmless
while it actually grants valuable authority?
\end{lstlisting}

\subsection{Wrong Signer and Wrong Authority}

Many bugs come from verifying a valid signature from the wrong
authority.

Examples:

\begin{itemize}
\tightlist
\item
  Recover signer but do not check signer equals
  \passthrough{\lstinline!owner!}.
\item
  Check token owner, but not approved operator.
\item
  Check message signer, but action affects another account.
\item
  Check guardian signature for owner-only action.
\item
  Accept a signature from any multisig owner instead of threshold
  result.
\item
  Treat relayer signature as user authorization.
\item
  Treat backend signature as proof that user consented.
\item
  Accept a contract wallet signature by checking ECDSA instead of
  ERC-1271.
\end{itemize}

Signature verification answers:

\begin{lstlisting}
This key signed this message.
\end{lstlisting}

Authorization must answer:

\begin{lstlisting}
Is this key allowed to authorize this transition for this asset in this state?
\end{lstlisting}

\subsection{Stale Authority}

An authorization may be valid when signed but unsafe when submitted.

Examples:

\begin{itemize}
\tightlist
\item
  A permit signed before a user revoked trust.
\item
  A limit order signed before market conditions changed.
\item
  A delegation signature signed before governance parameters changed.
\item
  A withdrawal signature signed before an account was frozen.
\item
  A session key valid after a user changed devices.
\item
  A contract-wallet signature that remains accepted after owner
  rotation.
\end{itemize}

Defenses:

\begin{itemize}
\tightlist
\item
  Short deadlines.
\item
  Nonce invalidation.
\item
  Versioned domains.
\item
  State-bound signatures.
\item
  Revocation registries.
\item
  Cancel functions.
\item
  Fill accounting.
\item
  Expected state checks.
\end{itemize}

Question:

\begin{lstlisting}
If this signature is submitted later, what assumptions may no longer hold?
\end{lstlisting}

\subsection{Key Compromise}

Key compromise means the attacker can generate valid signatures.

Once this happens, the chain cannot usually distinguish attacker
signatures from legitimate signatures. The signatures are valid.

Response depends on the authority:

\begin{itemize}
\tightlist
\item
  User account: sweep assets, revoke approvals, move NFTs, rotate
  smart-account owners if possible.
\item
  Admin key: pause protocol, rotate roles, inspect pending privileged
  actions.
\item
  Multisig signer: remove signer if threshold still safe.
\item
  Threshold share: refresh shares or rotate key if protocol supports it.
\item
  Validator key: stop duties, prevent slashable messages, rotate if
  supported.
\item
  Oracle key: remove publisher, invalidate bad reports, pause dependent
  markets.
\item
  Bridge key: pause bridge, rotate validator set, audit minted supply.
\item
  Deployment key: verify deployments and revoke CI/CD secrets.
\end{itemize}

Key-compromise analysis:

\begin{lstlisting}
What can the key sign?
What already-signed messages are pending?
What approvals or delegations remain active?
Can authority be revoked?
Can funds be moved before attacker moves them?
Can users observe compromise before damage?
Can protocol state distinguish pre-compromise and post-compromise signatures?
\end{lstlisting}

If the system has no key-compromise plan, the key is a single point of
final failure.

\subsection{Authorization Review Checklist}

For any signature-gated action, write:

\begin{lstlisting}
Action:
  What state transition is being authorized?

Signer:
  Which key, account, contract, role, or threshold policy authorizes it?

Message:
  What exact bytes are signed?

Domain:
  Chain, contract, version, function, module, source/destination, and operation type.

Replay:
  Nonce, sequence, nullifier, order ID, fill accounting, or deadline.

Verification:
  ECDSA, EdDSA, BLS, Schnorr, ERC-1271, multisig, threshold, or custom validation.

Signer-set policy:
  One key, m-of-n, threshold, validator set, oracle committee, or governance rule.

State dependency:
  Which state must still be true when the signature is submitted?

Failure mode:
  What happens if the key is stolen, lost, stale, malicious, or unavailable?

User comprehension:
  Can the signer understand the authority being granted?
\end{lstlisting}

This checklist is deliberately broader than cryptography because
authorization is broader than cryptography.

\section{Worked Example: Permit as Delegated Authority}

Consider an ERC-20 token implementing ERC-2612
\passthrough{\lstinline!permit!}:

\begin{lstlisting}
permit(owner, spender, value, deadline, v, r, s)
\end{lstlisting}

A user signs a typed message authorizing:

\begin{lstlisting}
owner = Alice
spender = Router
value = 10,000 USDC
deadline = Friday 17:00 UTC
nonce = 42
token = USDC
chain = Ethereum mainnet
verifying contract = USDC token contract
\end{lstlisting}

If verification succeeds, the token updates:

\begin{lstlisting}
allowance[Alice][Router] = 10,000 USDC
nonce[Alice] += 1
\end{lstlisting}

Now analyze it as authorization:

\begin{lstlisting}
Asset:
  Alice's token balance.

Authority granted:
  Router can transfer up to 10,000 USDC from Alice using transferFrom.

Signer:
  Alice's account.

Replay protection:
  Alice's permit nonce plus deadline.

Domain:
  token contract, chain ID, domain separator, typed Permit structure.

State transition:
  allowance changes; no tokens move yet.
\end{lstlisting}

The most common user misunderstanding:

\begin{lstlisting}
Alice thinks she signed a swap.
Actually she signed an approval.
\end{lstlisting}

The approval may be safe if:

\begin{itemize}
\tightlist
\item
  Router is the intended contract.
\item
  Amount is bounded.
\item
  Deadline is short.
\item
  Token contract consumes nonce correctly.
\item
  UI displays spender, amount, token, deadline, and chain.
\item
  The spender's later transfer is constrained by the swap logic Alice
  intended.
\end{itemize}

It may be unsafe if:

\begin{itemize}
\tightlist
\item
  Spender is malicious.
\item
  Amount is unlimited.
\item
  Deadline is far in the future.
\item
  Signature can be replayed.
\item
  UI hides that this is approval authority.
\item
  The signed domain is misleading.
\end{itemize}

The signature is not the end of the analysis. It is the beginning of the
authority trail.

\section{Practical Discipline: Translate Signatures Into Authority}

When you see a signature, translate it into this sentence:

\begin{lstlisting}
This signature allows [actor/verifier] to cause [state transition]
over [asset/state] in [domain], until [expiry],
provided [replay/state checks] pass.
\end{lstlisting}

Examples:

\begin{lstlisting}
This transaction signature allows Ethereum validators to include a transaction
from Alice's EOA that calls Router.swap on mainnet with nonce 17 and specified gas fields.
\end{lstlisting}

\begin{lstlisting}
This permit signature allows the token contract to set Router's allowance
over Alice's USDC to 10,000 until the deadline, if Alice's permit nonce is unused.
\end{lstlisting}

\begin{lstlisting}
This BLS signature allows the consensus protocol to count validator V's attestation
for this slot, fork, domain, and beacon block root.
\end{lstlisting}

If you cannot fill in that sentence, the authorization model is
incomplete.

\section{Review Questions}

\begin{enumerate}
\def\labelenumi{\arabic{enumi}.}
\tightlist
\item
  What does a digital signature prove, and what does it not prove?
\item
  Why is signature verification only one part of authorization?
\item
  What fields should a signed payload bind to reduce replay and confused-signing risk?
\item
  Why do transaction signing and message signing create different review questions?
\item
  How can a valid signature authorize an unsafe state transition?
\item
  What is the difference between an on-chain multisig, threshold signature scheme, MPC signing process, and hardware signer?
\item
  Why should private keys, validator keys, oracle keys, bridge keys, and admin keys be treated as authority-bearing assets?
\item
  Why does typed data improve signing safety without proving that a user understood the authorization?
\end{enumerate}

\section{Applied Exercises}

\subsection{Exercise 1: Signed-Payload Anatomy}

Given a signed approval or permit-like message, write an authorization-flow map with:

\begin{itemize}
\tightlist
\item
  signer;
\item
  verifying domain;
\item
  chain ID or domain identifier;
\item
  verifying contract or program;
\item
  nonce or sequence number;
\item
  deadline or expiry;
\item
  action being authorized;
\item
  state transition that becomes possible if the signature is accepted.
\end{itemize}

Expected output: one table and one paragraph explaining whether the signed payload is scoped narrowly enough.

\subsection{Exercise 2: Replay and Domain-Separation Review}

A protocol accepts signatures authorizing withdrawals. The message contains:

\begin{lstlisting}
withdraw(user, token, amount)
\end{lstlisting}

It omits chain ID, verifying contract, nonce, and deadline.

Identify at least four replay or confused-signing risks. Then rewrite the payload fields so the authorization is scoped to one protocol, one domain, one action, and one time window.

Expected output: a risk list and a revised message schema.

\subsection{Exercise 3: Compare Signing Arrangements}

Compare the following for a protocol treasury:

\begin{itemize}
\tightlist
\item
  one EOA;
\item
  on-chain multisig;
\item
  threshold signature or MPC policy;
\item
  hardware-backed signing process.
\end{itemize}

For each, list the authority it grants, the failure mode, the observability of misuse, and the recovery path.

Expected output: one comparison table.

\subsection{Exercise 4: Authorization Finding Draft}

Counterexample:

\begin{lstlisting}
A permit signature is valid forever.
The protocol does not include a nonce in the signed message.
The same signature can be submitted multiple times after the user's balance changes.
\end{lstlisting}

Write a concise finding with:

\begin{itemize}
\tightlist
\item
  root cause;
\item
  affected authority;
\item
  exploit path;
\item
  impact;
\item
  recommended mitigation;
\item
  regression test.
\end{itemize}

Expected output: a short review note, not a full report.

\section{Key Takeaways}

\begin{itemize}
\tightlist
\item
  Private keys are authority material. Evaluate them by what transitions
  they can authorize.
\item
  Public keys verify signatures; addresses identify accounts, scripts,
  contracts, programs, or derived public-key material.
\item
  Seed phrases and HD wallet roots can control many accounts across many
  chains.
\item
  A wallet is an interface to authority, not the same thing as the
  account.
\item
  Signature algorithms rely on assumptions about private-key secrecy,
  randomness, encoding, domain separation, public-key validation, and
  implementation correctness.
\item
  ECDSA nonce failures can reveal private keys; ECDSA malleability can
  break naive replay protection.
\item
  EdDSA, BLS, and Schnorr have different strengths and different
  integration risks.
\item
  BLS aggregation and threshold signatures are powerful but require
  explicit signer-set, proof-of-possession, nonce, and domain handling.
\item
  Message signing is often more dangerous than transaction signing
  because the meaning is defined by the verifier.
\item
  EIP-712 improves structured signing but does not remove the need for
  nonces, deadlines, and correct domain separation.
\item
  Multisig, threshold signatures, MPC, and hardware wallets reduce some
  risks while introducing others.
\item
  Authorization means a valid signature plus a correct policy decision
  about what transition that signer may cause.
\end{itemize}

\section{Further Reading}

\begin{itemize}
\tightlist
\item
  NIST,
  \href{https://nvlpubs.nist.gov/nistpubs/FIPS/NIST.FIPS.186-5.pdf}{``FIPS
  186-5: Digital Signature Standard''}. Primary reference for digital
  signature concepts, ECDSA, EdDSA, and key-protection warnings.
\item
  ethereum.org,
  \href{https://ethereum.org/developers/docs/accounts/}{``Ethereum
  Accounts''}. Useful distinction between EOAs, contract accounts,
  private keys, addresses, and wallets.
\item
  Pieter Wuille,
  \href{https://github.com/bitcoin/bips/blob/master/bip-0032.mediawiki}{``BIP-32:
  Hierarchical Deterministic Wallets''}. HD wallet derivation, extended
  keys, hardened derivation, and wallet tree structure.
\item
  Marek Palatinus et al.,
  \href{https://github.com/bitcoin/bips/blob/master/bip-0039.mediawiki}{``BIP-39:
  Mnemonic code for generating deterministic keys''}. Mnemonic-to-seed
  construction and seed phrase framing.
\item
  Josefsson and Liusvaara,
  \href{https://datatracker.ietf.org/doc/html/rfc8032}{``RFC 8032:
  Edwards-Curve Digital Signature Algorithm''}. EdDSA standard
  reference.
\item
  Pieter Wuille et al.,
  \href{https://github.com/bitcoin/bips/blob/master/bip-0340.mediawiki}{``BIP-340:
  Schnorr Signatures for secp256k1''}. Bitcoin Schnorr signature
  standard.
\item
  Boneh et al.,
  \href{https://datatracker.ietf.org/doc/html/draft-irtf-cfrg-bls-signature}{``BLS
  Signatures''}. CFRG Internet-Draft for BLS signatures and aggregation
  considerations; check current status before treating it as a finalized
  standard.
\item
  Ethereum Consensus Specs,
  \href{https://ethereum.github.io/consensus-specs/specs/phase0/beacon-chain/\#bls-signatures}{``Beacon
  Chain: BLS Signatures''}. Ethereum consensus-layer context for BLS
  usage.
\item
  Connolly et al.,
  \href{https://www.rfc-editor.org/rfc/rfc9591.html}{``RFC 9591:
  FROST''}. Threshold Schnorr signing protocol and security
  considerations.
\item
  Ethereum Improvement Proposals,
  \href{https://eips.ethereum.org/EIPS/eip-191}{``EIP-191: Signed Data
  Standard''}. Ethereum signed-data prefixing and versioning.
\item
  Ethereum Improvement Proposals,
  \href{https://eips.ethereum.org/EIPS/eip-712}{``EIP-712: Typed
  structured data hashing and signing''}. Structured data signing and
  domain separators.
\item
  Ethereum Improvement Proposals,
  \href{https://eips.ethereum.org/EIPS/eip-155}{``EIP-155: Simple Replay
  Attack Protection''}. Chain ID replay protection for Ethereum
  transactions.
\item
  Ethereum Improvement Proposals,
  \href{https://eips.ethereum.org/EIPS/eip-2}{``EIP-2: Homestead
  Hard-fork Changes''}. Includes Ethereum's high-s transaction signature
  rule.
\item
  Ethereum Improvement Proposals,
  \href{https://eips.ethereum.org/EIPS/eip-2612}{``ERC-2612: Permit
  Extension for EIP-20 Signed Approvals''}. Permit-based signed token
  approvals.
\item
  Ethereum Improvement Proposals,
  \href{https://eips.ethereum.org/EIPS/eip-1271}{``ERC-1271: Standard
  Signature Validation Method for Contracts''}. Signature validation for
  contract accounts.
\item
  Safe, \href{https://docs.safe.global/home/what-is-safe}{``What is
  Safe?''}. Practical reference for smart-account and multisig-style
  custody.
\item
  NIST,
  \href{https://nvlpubs.nist.gov/nistpubs/SpecialPublications/NIST.SP.800-57pt1r5.pdf}{``SP
  800-57 Part 1 Revision 5: Recommendation for Key Management''}.
  General key-management lifecycle and cryptographic-key guidance.
\item
  NIST,
  \href{https://nvlpubs.nist.gov/nistpubs/FIPS/NIST.FIPS.140-3.pdf}{``FIPS
  140-3: Security Requirements for Cryptographic Modules''}. Hardware
  and software cryptographic-module security requirements.
\end{itemize}
